% =============================================================
%  Journal of AI Ethics (JOAIE) — Official Paper Template
%  Version 1.0 | April 2026
%  https://journalofaiethics.org
% =============================================================
%
%  INSTRUCTIONS
%  ------------
%  1. Fill in the metadata fields below (title, author, date).
%  2. Write your paper in the sections provided.
%  3. Leave the Reference Code field blank on first submission.
%     The editor will assign a code (JOAIE-YYYY-XXX) on acceptance.
%  4. Generate HTML and PDF outputs with Pandoc before submitting:
%
%       pandoc JOAIE-YYYY-XXX.tex -o JOAIE-YYYY-XXX.html \
%         --standalone --css=../../styles.css
%
%       pandoc JOAIE-YYYY-XXX.tex -o JOAIE-YYYY-XXX.pdf
%
%  5. Place all three files (.tex, .html, .pdf) in a folder named
%     papers/JOAIE-YYYY-XXX/ and open a pull request.
%
%  REQUIREMENTS SUMMARY
%  --------------------
%  - Abstract: maximum 250 words
%  - Keywords: 3–8 comma-separated terms
%  - References: APA 7th edition (use natbib / apalike)
%  - All sections below are required unless marked [optional]
%  - Disclose any use of AI tools in a note at the end of the paper
% =============================================================

\documentclass[12pt, a4paper]{article}

% --- Layout ---
\usepackage[margin=1in]{geometry}
\usepackage{setspace}
\onehalfspacing

% --- Typography ---
\usepackage[T1]{fontenc}
\usepackage[utf8]{inputenc}
\usepackage{microtype}          % Improved text justification

% --- Links & URLs ---
\usepackage[hidelinks]{hyperref}
\usepackage{url}

% --- References (APA-style) ---
\usepackage{natbib}
\bibliographystyle{apalike}

% --- Tables & Figures ---
\usepackage{booktabs}           % Professional tables
\usepackage{graphicx}           % Include images
\usepackage{float}              % Better figure placement
\usepackage{caption}
\captionsetup{font=small, labelfont=bf, skip=6pt}

% --- Mathematics [optional] ---
\usepackage{amsmath}
\usepackage{amssymb}

% --- Code listings [optional] ---
\usepackage{listings}
\lstset{basicstyle=\small\ttfamily, breaklines=true}

% =============================================================
%  METADATA — Fill in before submission
% =============================================================

\title{Your Paper Title}

\author{
  Your Name \\
  \small Your Institution or "Independent Researcher" \\
  \small \href{mailto:you@example.com}{you@example.com}
}

\date{Month Year}

% --- Reference Code (leave blank on first submission) ---
% JOAIE-YYYY-XXX

% =============================================================
%  DOCUMENT
% =============================================================

\begin{document}
\maketitle

% --- Horizontal rule beneath title ---
\noindent\rule{\linewidth}{0.4pt}
\vspace{0.5em}

% =============================================================
%  ABSTRACT
%  Maximum 250 words. No citations in the abstract.
% =============================================================

\begin{abstract}
\noindent
Your abstract here. Summarise the research question, methodology,
key findings, and implications. Maximum 250 words. Do not include
citations in the abstract.
\end{abstract}

\noindent\textbf{Keywords:} keyword1, keyword2, keyword3, keyword4

\noindent\rule{\linewidth}{0.4pt}
\vspace{1em}

% =============================================================
%  SECTIONS
%  All sections are required. You may add subsections as needed.
% =============================================================

\section{Introduction \& Research Question}
%
% State the problem or phenomenon you are addressing, explain why
% it matters, and close with a clear, specific research question.
% Typical length: 400–600 words.

\section{Background \& Literature Review}
%
% Situate your work in the existing literature. What has been
% established, and where does your contribution fit? Cite
% relevant work using \citep{} and \citet{} commands.
% Typical length: 600–900 words.

\section{Methodology}
%
% Explain how you conducted the research or constructed the
% argument. For empirical work: describe your data, experimental
% design, and analysis approach. For theoretical work: describe
% the analytical framework and how you applied it.
% Typical length: 500–800 words.

\section{Findings \& Analysis}
%
% Present your results and analyse what they mean. For empirical
% work: report your data with appropriate statistical context.
% For theoretical work: develop your argument step by step.
% Typical length: 800–1200 words.

\section{Discussion}
%
% Interpret your findings in light of the broader literature.
% Address limitations of your study. Discuss implications for
% practice, policy, or further research.
% Typical length: 400–600 words.

\section{Conclusion}
%
% Summarise the key contribution and what should be taken away.
% No new information here — only synthesis. End with a clear
% statement of the paper's normative or practical implication.
% Typical length: 200–350 words.

% =============================================================
%  AI TOOL DISCLOSURE [required if applicable]
% =============================================================

% Uncomment and fill in if you used AI tools in this research:
%
% \section*{AI Tool Disclosure}
% This paper was prepared [with / without] the assistance of AI
% writing tools. [Describe what tool was used, for what purpose,
% and how outputs were reviewed or verified.]

% =============================================================
%  REFERENCES
%  Uses natbib with apalike style (APA 7th edition approximation)
%  Store your .bib file as references.bib in the same folder.
% =============================================================

\bibliography{references}

\end{document}
